% Optional structural continuation; all formulas are exact.
% Reuses the Padé polynomials pi_n, constants c_n and h_n from moment_rank.tex.
\section{A fixed-size determinant for every pair of even orders}
\label{sec:general-even-reduction}

The preceding argument has a useful extension even where the final
nonvanishing problem remains unsolved.  For $p=2a$, $q=2b$, with
$1\leq a\leq b$, the optimal-prefix matrix is an explicitly scaled
rank-$a$ perturbation of a shifted Bernoulli Hankel matrix.  The Hankel
matrix is always nonsingular.  Moreover, the remaining determinant
has size $a$, independently of $q$.

For clarity, the general inherited moment notation is
\[
 D_{p,q}(s)=\sum_{n\ge0}A_p(n)A_q(n)(n+1/2)^{-s},
 \qquad \re s>3-p-q,
\]
continued meromorphically thereafter, and $D_r=D_{p,q}(-2r)$.
The same centered-parity argument proves regularity at these even
arguments. The forward and inverse Bernoulli transforms
\eqref{eq:forward-foundation}--\eqref{eq:inverse-foundation} remain valid,
now with $C_j$ the constant term at $N\to\infty$ of
\begin{align*}
 &N^{2j+1}A_p(N)A_q(N)
 +\sum_{k=1}^N k^{2j+1-q}A_p(k)\\
 &\hspace{15mm}
 +\sum_{k=1}^N k^{2j+1-p}A_q(k)
 +\sum_{k=1}^N k^{2j+1-p-q}.
\end{align*}
For $j\ge b$ the inherited formula is
$C_j=\sum_{k=0}^{b-1}a_{jk}Z_k$, with $Z_0=-1/2$ and
$Z_k=\zeta(2k)$ for $k\ge1$. Here the coefficients depend on
the indicated pair $(p,q)$. With $d=2j+1$ and
$e=d-p-q+2k+1$, they are zero for $e<0$, and otherwise
\begin{equation}\label{eq:general-source-entries}
 a_{jk}=B_e\left\{
 \frac{\binom{d-q+1}{p-2k}}{d-q+1}
 +\frac{\binom{d-p+1}{q-2k}}{d-p+1}\right\}.
\end{equation}
Each summand is omitted unless its lower binomial index lies between
$2$ and its upper index. These are the exact general formulas of
\cite[Section 7]{CanonicalFifth}; the finite summation argument of
Section~\ref{sec:tail-foundation} has the same form. In the formal
coordinate order $C_0,\ldots,C_{b-1},Z_0,\ldots,Z_{b-1}$, the
complete first-$q$ determinant is $\det M_{a,b}/(4b-1)!!$.
The new theorems below evaluate or reduce the remaining square minor.

Put
\[
 \mu_j=\begin{cases}B_{2j}/(2j)!,&j\geq0,\\0,&j<0,\end{cases}
 \qquad H_{a,b}=(\mu_{t+k-a+1})_{0\leq t,k<b}.
\]
The convention $\mu_j=0$ for $j<0$ is part of this definition; it is
not an analytic continuation of the formula for positive indices.

\begin{lemma}[Factorization of the general source matrix]
\label{lem:general-even-factor}
Let $M_{a,b}=(a_{b+t,k})_{0\leq t,k<b}$, with the source's convention
that a binomial summand is omitted when its lower index is less than
$2$.  Define
\begin{align*}
 R_{tt}&=(2b+2t-2a+1)!,& E_{kk}&=\frac1{(2b-2k)!},\\
 U_{tk}&=\mathbf1_{k<a}\,\mu_{t+k-a+1}
 \frac{(2t+1)!(2b-2k)!}
 {(2a-2k)!(2b+2t-2a+1)!}.
\end{align*}
Then
\[
 M_{a,b}=R(H_{a,b}+U)E,
 \qquad \operatorname{rank}U=a.
\]
\end{lemma}

\begin{proof}
In the source entry formula put $j=b+t$ and
$e=2(t+k-a+1)$.  When $e\geq0$, cancellation of the factorial in
$B_e$ gives
\[
 a_{b+t,k}=\mu_{t+k-a+1}\left\{
 \frac{(2b+2t-2a+1)!}{(2b-2k)!}
 +\mathbf1_{k<a}\frac{(2t+1)!}{(2a-2k)!}\right\}.
\]
When $e<0$, both sides vanish by the stated convention.  This proves
the factorization.  Only the first $a$ columns of $U$ are nonzero.
Their first $a$ rows form an antitriangular matrix with nonzero
antidiagonal, because $\mu_{t+k-a+1}=0$ for $t+k<a-1$ and
$\mu_0=1$.  Thus its rank is exactly $a$.
\end{proof}

\begin{theorem}[Nonsingularity and inertia of all the shifted matrices]
\label{thm:shifted-H-inertia}
The real symmetric matrix $H_{a,b}$ is nonsingular for every
$1\leq a\leq b$.  Its numbers of positive and negative eigenvalues
are, respectively,
\[
 \begin{cases}
 (b-h,h),&a=2h,\\
 (h+1,b-h-1),&a=2h+1.
 \end{cases}
\]
Consequently
\[
 \operatorname{sgn}\det H_{a,b}=
 \begin{cases}(-1)^h,&a=2h,\\
 (-1)^{b-h-1},&a=2h+1.
 \end{cases}
\]
\end{theorem}

\begin{proof}
Euler's even-zeta evaluation gives
\[
 \mu_r=(-1)^{r-1}\,2\sum_{m\geq1}x_m^r,
 \qquad r\geq1,\qquad x_m=(2\pi m)^{-2}.
\]
Conjugate $H_{a,b}$ by $\operatorname{diag}(1,-1,1,\ldots)$ and
call the result $G$.  For $t+k\geq a$ its entries are
$(-1)^a2\sum_{m\geq1}x_m^{t+k-a+1}$; on $t+k=a-1$ they are
$(-1)^{a-1}$, and below that antidiagonal they are zero.
Every moment matrix
$\bigl(2\sum_m x_m^{i+j+s}\bigr)_{0\leq i,j<d}$ with $s\geq1$
is positive definite: its quadratic form is
$2\sum_m x_m^s P(x_m)^2$, and a nonzero polynomial cannot vanish
at all the distinct positive points $x_m$.

First let $a=2h$.  Partition after the first $h$ indices:
\[
 G=\begin{pmatrix}0&B\\B^{\mathsf T}&P\end{pmatrix}.
\]
Here $P$ is positive definite.  The first $h$ columns of $B$ form
an antitriangular square matrix with antidiagonal $-1$, so $B$ has
row rank $h$.  The Schur complement
$-BP^{-1}B^{\mathsf T}$ is negative definite.  Congruence elimination
therefore gives $b-h$ positive and $h$ negative eigenvalues.

Now let $a=2h+1$ and partition after the first $h+1$ indices:
\[
 G=\begin{pmatrix}ee^{\mathsf T}&B\\
 B^{\mathsf T}&-P\end{pmatrix},
 \qquad e=(0,\ldots,0,1)^{\mathsf T}.
\]
The matrix $P$ is positive definite.  In the first $h$ rows of $B$,
its first $h$ columns form an antitriangular matrix with
antidiagonal $1$.  Hence
$ee^{\mathsf T}+BP^{-1}B^{\mathsf T}$ is positive definite:
if its quadratic form vanishes, first $e^{\mathsf T}v=0$ and then
the indicated minor forces the remaining $h$ components of $v$ to
vanish.  Congruence elimination gives $h+1$ positive and
$b-h-1$ negative eigenvalues.  The case $a=b=1$ is simply
$H_{1,1}=(1)$; empty blocks in the argument may otherwise be
omitted.  Diagonal sign conjugation preserves inertia.
\end{proof}

There is a more explicit reduction than a matrix determinant lemma
with an unevaluated $H_{a,b}^{-1}$.  Retain the monic polynomials
$\pi_n$ and positive numbers $h_n,c_n$ constructed in the preceding
Pad\'e calculation.  Write
\[
 \pi_n(x)=\sum_{t=0}^n p_{nt}x^t,
 \qquad
 p_{nt}=\frac{2(2n+2t+1)!(2n+1)!}
 {(4n+2)!(2n-2t+1)!(2t)!}.
\]
For $0\leq i,k<a$ set $n_i=b-a+i$, and define the $a$-square matrix
\begin{equation}\label{eq:general-small-A}
 A_{ik}=\sum_{t=0}^{n_i}p_{n_i t}\mu_{t+k-a+1}
 \left\{1+
 \frac{(2t+1)!(2b-2k)!}
 {(2a-2k)!(2b+2t-2a+1)!}\right\}.
\end{equation}
These are finite rational sums with explicitly specified coefficients.

\begin{theorem}[Reduction to a determinant of size $a$]
\label{thm:general-small-determinant}
For $1\leq a\leq b$,
\begin{equation}\label{eq:general-small-det}
 \det M_{a,b}=
 (-1)^{a(b-a)}
 \frac{\prod_{t=0}^{b-1}(2b+2t-2a+1)!}
      {\prod_{k=0}^{b-1}(2b-2k)!}
 \left(\prod_{n=0}^{b-a-1}h_n\right)\det A.
\end{equation}
All displayed scalar factors are nonzero.  Thus the remaining
optimal-prefix nonvanishing question is exactly $\det A\ne0$.
This reduction by itself does not assert that nonvanishing.
\end{theorem}

\begin{proof}
Let $\mathcal L_a(x^r)=\mu_{r-a+1}$ and recall
$\mathcal N(x^r)=\mu_{r+1}$.  Then
$\mathcal L_a(x^ag)=\mathcal N(g)$.  Replace row $n$ of
$H_{a,b}+U$ by the linear combination whose coefficients are those
of $\pi_n$, simultaneously for $0\leq n<b$.  This is a triangular
row transformation with diagonal $1$.

For a column $k\geq a$ the new entry is
\[
 \mathcal L_a(x^k\pi_n)=\mathcal N(x^{k-a}\pi_n).
\]
It vanishes when $k-a<n$, and equals $h_n$ when $k-a=n$.
The last $a$ rows consequently vanish in the last $b-a$ columns,
and their first $a$ columns are exactly $A$.  The first $b-a$ rows
in the last $b-a$ columns form an upper triangular matrix with
diagonal $h_0,\ldots,h_{b-a-1}$.  Moving the first $a$ columns past
the other $b-a$ columns gives the sign $(-1)^{a(b-a)}$.
Taking determinants and then restoring $R$ and $E$ proves the formula.
\end{proof}

For reference, the same reduction gives a fully explicit
small determinant for the unperturbed Hankel matrix.  Put
\[
 f_\ell(n)=\frac1{(2\ell+1)!}
 \prod_{r=0}^{2\ell-1}(2n-r)(2n+3+r),\qquad f_0(n)=1.
\]
Then
\begin{equation}\label{eq:general-H-small-det}
 \det H_{a,b}=(-1)^{a(b-a)}
 \left(\prod_{n=0}^{b-a-1}h_n\right)
 \left(\prod_{i=0}^{a-1}c_{n_i}\right)
 \det\bigl(f_{a-1-k}(n_i)\bigr)_{0\leq i,k<a}.
\end{equation}
Indeed,
\[
 \mathcal L_a(x^k\pi_n)
 =[z^{n+k-a+1}]F(z)q_n(z)
 =[z^{n+k-a+1}]r_n(z)
 =c_n f_{a-1-k}(n),
\]
where a negative coefficient index means zero.  The factorial
coefficient formula for $r_n$ proves the last equality, including
the zero when $n<a-1-k$.  In particular, the inertia theorem gives
the exact sign statement
\[
 (-1)^{a(a-1)/2}
 \det\bigl(f_{a-1-k}(b-a+i)\bigr)_{0\leq i,k<a}>0.
\]

\paragraph{A concrete route for the remaining even orders.}
For each fixed $a$, the entries in \eqref{eq:general-small-A} can be
reduced to finitely many nearby Bernoulli numbers by weighted
binomial Bernoulli identities.  More explicitly, if
$n=b-a+i$, $\ell=a-1-k$, and $N=2(n-\ell)+1\geq1$, the perturbation
part of $A_{ik}$ is
\[
 \frac{2(2n+1)!(2b-2k)!}
 {(4n+2)!(2a-2k)!N!}
 \sum_{r=0}^{n-\ell}\binom N{2r}B_{2r}
 (2r+2\ell+1)
 (2n+2\ell+2r+1)^{\underline{2i}}.
\]
The underline denotes a falling factorial.  The polynomial weight
has degree $2i+1\leq2a-1$.  Expressing it in the basis
$(N-2r)^{\underline d}$ and using
\[
 \sum_{j=0}^N\binom Nj B_j(N-j)^{\underline d}
   =N^{\underline d}B_{N-d}(1)
\]
therefore reduces it to Bernoulli numbers of orders at most $2a-1$
below $N$, with the omitted $j=1$ term restored explicitly.
When $n<\ell$ the perturbation entry is zero.  The small cases
$N=1,3,\ldots,2i+1$ require retaining $B_1(1)=1/2$ and $B_0=1$;
discarding them as generic odd-index zeros would be an error.
This is a finite algebraic program for each $p$, followed by a
genuine nonvanishing problem.  The sign theorem for $H_{a,b}$
does not determine the sign or nonvanishing of its perturbation.
